\section{Algorithm and Implementation Details}

The project uses a ResNet18 image classifier \citep{he2016resnet} with a squeeze-and-excitation style channel attention block \citep{hu2018senet}. Training uses pretrained torchvision weights and a task-specific seven-class head. The implementation is written in PyTorch \citep{paszke2019pytorch} for single-image dermoscopic classification.

\begin{table}[H]
    \centering
    \caption{Current algorithm summary.}
    \label{tab:algorithm}
    \begin{tabular}{>{\raggedright\arraybackslash}p{0.28\textwidth} >{\raggedright\arraybackslash}p{0.62\textwidth}}
        \toprule
        Component & Current implementation \\
        \midrule
        Task type & Seven-class classification \\
        Input & One RGB dermoscopic image \\
        Backbone & ResNet18 \\
        Training strategy & Transfer learning \\
        Attention module & SE-style channel attention \\
        Classifier head & 512$\rightarrow$256$\rightarrow$7 with ReLU and Dropout(0.4) \\
        Image preprocessing & Resize \(224\times224\), tensor, ImageNet normalization \\
        Training augmentation & H/V flips, 20$^\circ$ rotation, mild color jitter \\
        Loss function & Cross-entropy \\
        Optimizer & Adam \\
        Learning-rate schedule & Cosine annealing, minimum \(10^{-6}\) \\
        Batch size and epochs & 32, 30 \\
        Inference rule & Top-1 softmax with confidence \\
        Extra deployment logic & Blank-image heuristic, threshold 0.5 \\
        Main strength & Compact and deployable \\
        Main limitation & No validation stage, class imbalance \\
        \bottomrule
    \end{tabular}
\end{table}

This design balances standard transfer learning with lightweight deployment. Its main limits are unchanged: the model predicts one label per image, does not localize lesions, and returns confidence scores that are not formally calibrated.
